\chapter{Faire ou acheter (build against buy)}\label{ch:fm:build-against-buy}

En novembre 2018, Virtu Financial accepte de racheter Investment Technology Group, un courtier à commission, et en mars 2019 elle finalise le rachat pour environ
1.0~milliard de dollars en numéraire, soit 30.30~\$ par action. Ses motifs affichés étaient d'offrir à ses clients une gamme complète de services d'exécution, «~y compris
la technologie de trading transparent et de workflow, les analyses et les solutions de liquidité~», fonctionnant sur l'infrastructure propre de Virtu, et environ 123~millions de dollars par an d'économies de
charges dans les deux ans. Une société peut écrire le logiciel dont elle a besoin, le licencier, ou racheter l'entreprise qui l'a écrit avec ses clients
et ses équipes. La dernière option est la plus coûteuse pour acquérir un logiciel, et parfois la seule pour acquérir une activité bâtie dessus.

\section{Ce que les sociétés construisent et ce qu'elles achètent}

\begin{definition}[Faire ou acheter]\label{def:fm:build-against-buy:bab}
La décision \emph{faire ou acheter}\index{faire ou acheter} (build against buy) consiste à choisir, pour une capacité dont la société a besoin, entre la développer avec ses propres
ingénieurs, en prenant une licence ou un abonnement auprès d'un fournisseur, ou en combinant les deux approches (un noyau fourni par un fournisseur, complété par les extensions propres à la société), sur
toute la durée de vie de la capacité.
\end{definition}

La règle habituelle veut qu'une société construise ce qui la distingue et achète ce qui la rend identique aux autres. Le moteur de cotation et les contrôles de risque d'un teneur de marché sont
son edge, qu'elle construit (les livres~11 et~13)~; sa paie, sa messagerie et une bonne part de son traitement post-négociation sont identiques à ceux de tout le monde, et elle
les achète. Les cas difficiles se situent entre les deux~: les systèmes de gestion des ordres et de l'exécution, OMS et EMS (livre~10), le référentiel des titres (livre~7), la plateforme d'apprentissage
automatique (livre~12). Tout le monde en a besoin, ils touchent à l'edge de la société par ses bords, et ils coûtent cher dans les deux sens.

\begin{definition}[Coût total de possession]\label{def:fm:build-against-buy:tco}
Le \emph{coût total de possession}\index{coût total de possession} (total cost of ownership) d'une capacité est la valeur actuelle de chaque coût lié à sa possession sur toute sa durée de vie~:
les frais de développement ou de licence, d'intégration et de personnalisation, d'infrastructure, de support et de maintenance, le temps des ingénieurs et ce qu'ils auraient
construit autrement, et le coût de sortie.
\end{definition}

\section{Coût total de possession}

L'exemple du chapitre est un système de gestion des ordres sur sept ans (le \cref{lst:fm:build-against-buy:costs}). Le construire en interne mobilise sept ingénieurs
pendant deux ans, puis trois par an pour le maintenir, plus 0.2~million de dollars par an d'infrastructure. L'achat coûte une licence de 1.2~millions de dollars qui augmente de 5~\% par
an, un coût d'intégration de 1.0~million de dollars et deux ingénieurs la première année, et un ingénieur de support chaque année. Une année-ingénieur coûte 350\,000~\$
toutes charges comprises, et la société ajoute 30~\% pour le travail que l'ingénieur aurait fait autrement~: 455\,000~\$. Le taux d'actualisation est de 8~\%. Tous les chiffres sont
illustratifs~; l'ouvrage de Boehm \emph{Software Engineering Economics} (1981) est la référence classique pour estimer les postes de développement et de maintenance.

\omcode{code/firm/buildbuy/firm_buildbuy.py}{41}{60}{Les deux chemins de coût et leur coût actualisé net~: les années de développement de la construction et sa queue de
maintenance, et la licence avec sa hausse annuelle, l'intégration et le support.}\label{lst:fm:build-against-buy:costs}

\begin{omfigure}
\begin{tikzpicture}
\begin{axis}[width=10.5cm,height=5cm,ybar,bar width=10pt,xlabel={\small année},ylabel={\small coût (millions de \$)},ymin=0,ymax=4,
  xtick={1,2,3,4,5,6,7},tick label style={font=\footnotesize},grid=major,grid style={gray!20},table/col sep=comma,
  legend cell align=left,legend style={font=\scriptsize,at={(0.5,-0.3)},anchor=north,/tikz/every even column/.append style={column sep=8pt}},legend columns=2]
\addplot[fill=omDef!55,draw=omDef] table[x=year,y=build]{figdata/desk/20-build-against-buy/paths.csv};
\addplot[fill=omProp!55,draw=omProp] table[x=year,y=buy]{figdata/desk/20-build-against-buy/paths.csv};
\legend{construction,achat}
\end{axis}
\end{tikzpicture}

\omcaption{Les coûts annuels des deux chemins~: les deux années de développement de la construction et sa queue de maintenance constante, et l'année d'intégration
de l'achat et sa licence croissante. Les années-ingénieur incluent le coût d'opportunité de 30~\%. Données~: \texttt{fm\_buildbuy.paths}.}\label{fig:fm:build-against-buy:paths}
\end{omfigure}

La construction coûte 3.19~millions de dollars chacune des deux premières années et 1.57~millions de dollars par an ensuite (la \cref{fig:fm:build-against-buy:paths})~; l'achat coûte 3.57~millions de dollars la première année et
ensuite 1.72~millions de dollars, qui augmente avec la licence jusqu'à 2.06~millions de dollars la septième année. Sans actualisation, la construction coûte 14.2~millions de dollars sur sept ans et l'achat
14.9~millions de dollars. L'actualisation favorise l'achat, puisque les coûts de la construction arrivent en premier~: les coûts actuels nets sont 11.04~millions de dollars et 11.30
millions de dollars (la \cref{fig:fm:build-against-buy:horizon}).

\begin{omfigure}
\begin{tikzpicture}
\begin{axis}[width=10.5cm,height=5.2cm,xlabel={\small horizon (années)},ylabel={\small coût actuel net (millions de \$)},xmin=1,xmax=12,ymin=0,ymax=18,
  tick label style={font=\footnotesize},grid=major,grid style={gray!20},table/col sep=comma,
  legend cell align=left,legend style={font=\scriptsize,at={(0.5,-0.3)},anchor=north,/tikz/every even column/.append style={column sep=8pt}},legend columns=2]
\addplot[omDef,thick,mark=*,mark size=1.2pt] table[x=years,y=build]{figdata/desk/20-build-against-buy/horizon.csv};
\addplot[omProp,thick,dashed,mark=square*,mark size=1.2pt] table[x=years,y=buy]{figdata/desk/20-build-against-buy/horizon.csv};
\draw[gray,dotted] (axis cs:7,0) -- (axis cs:7,18);
\legend{construction,achat}
\end{axis}
\end{tikzpicture}

\omcaption{Coût actuel net de la construction et de l'achat du système de gestion des ordres, pour des horizons d'un à douze ans. L'achat est moins cher de
l'année 2 à l'année 6~; à partir de l'année 7 la construction est moins chère et l'écart se creuse. Données~: \texttt{fm\_buildbuy.by\_horizon}.}\label{fig:fm:build-against-buy:horizon}
\end{omfigure}

\begin{proposition}[Le coût de l'ingénieur qui tranche]\label{prop:fm:build-against-buy:wage}
Si chaque année-ingénieur, sur l'un ou l'autre chemin, coûte $w(1+\kappa)$ et que tous les autres coûts ne dépendent pas de $w$, la différence entre les coûts actuels nets de
la construction et de l'achat est linéaire en $w$, $\Delta(w)=a+bw$, avec $b$ la valeur actuelle des années-ingénieur de la construction moins celles de l'achat, multipliée par
$1+\kappa$. Lorsque $b>0$, la construction est moins chère exactement lorsque $w<w^*=-a/b$.
\end{proposition}

\begin{proof}
Le coût de chaque année sur chaque chemin est une constante plus un nombre d'années-ingénieur multiplié par $w(1+\kappa)$~; l'actualisation et la sommation conservent la linéarité. La
construction utilise davantage d'années-ingénieur ($b>0$), donc $\Delta$ croît avec $w$ et change de signe une seule fois, en $-a/b$.
\end{proof}

Sur sept ans, le coût chargé de l'année-ingénieur au seuil d'indifférence est de 363\,000~\$~: au coût de 350\,000~\$ retenu par la société, la construction l'emporte de 0.26~million de dollars de coût actuel,
et une société qui paie plus cher ses ingénieurs devrait acheter. À partir de la septième année, la construction reste moins chère à tout horizon plus long. La marge est mince,
ce qui est le cas courant~: la décision se joue rarement sur le scénario de base et souvent sur ce qui pourrait mal tourner.

\section{Dépendance au fournisseur et option de changement}

\begin{definition}[Dépendance au fournisseur, coût de changement]\label{def:fm:build-against-buy:lockin}
La \emph{dépendance au fournisseur}\index{dépendance au fournisseur} (vendor lock-in) est la situation où l'on dépend d'un fournisseur au point que partir coûte cher~: données dans ses formats, flux de travail bâtis autour
de son produit, intégrations, compétences et contrats. Le \emph{coût de changement}\index{coût de changement} (switching cost) est ce que coûte le départ~: migration, fonctionnement en parallèle,
la formation des équipes, les frais de résiliation et le risque de perturbation.
\end{definition}

La dépendance permet à un fournisseur d'augmenter son prix au renouvellement, à hauteur du \omterm{def:fm:build-against-buy:lockin}{coût de changement}, sans perdre le client. Le chapitre modélise le risque par un arbre binomial~:
à chaque renouvellement, la licence peut augmenter de 30~\% avec une probabilité de 0.3~; changer de fournisseur coûte 2.0~millions de dollars et remet le prix à sa valeur initiale. Avec
le risque de prix, le coût actuel espéré de l'achat passe de 11.30~millions de dollars à 13.46~millions de dollars, soit 2.42~millions de dollars de plus que la construction. L'option de
changer de fournisseur, exercée quand elle rapporte, ramène ce coût à 12.90~millions de dollars. L'option vaut 0.56~million de dollars, et sa valeur justifie de garder
des \omterm{def:fm:build-against-buy:lockin}{coûts de changement} faibles (le \cref{lst:fm:build-against-buy:switch}). La meilleure politique est simple~: changer dès que le prix a monté deux fois, de 69~\%,
sauf la dernière année, où une seule année d'économies ne suffit plus à rentabiliser le changement (la \cref{fig:fm:build-against-buy:tree}).

\omcode{code/firm/buildbuy/firm_buildbuy.py}{79}{106}{L'option de changer de fournisseur~: induction à rebours sur un arbre des prix de renouvellement, avec
et sans la possibilité de changer de fournisseur.}\label{lst:fm:build-against-buy:switch}

\begin{omfigure}
\begin{tikzpicture}[font=\footnotesize,>=stealth,n/.style={draw,rounded corners,fill=omDef!12,minimum width=1.2cm,align=center},
  s/.style={draw,rounded corners,fill=omProp!30,minimum width=1.2cm,align=center}]
\node[n] (a) at (0,0) {1.00};
\node[n] (b1) at (3,1.1) {1.30}; \node[n] (b0) at (3,-1.1) {1.00};
\node[s] (c2) at (6,2.2) {1.69\\{\scriptsize changement}}; \node[n] (c1) at (6,0) {1.30}; \node[n] (c0) at (6,-2.2) {1.00};
\draw[->] (a) -- node[above,font=\scriptsize] {0.3} (b1); \draw[->] (a) -- node[below,font=\scriptsize] {0.7} (b0);
\draw[->] (b1) -- (c2); \draw[->] (b1) -- (c1); \draw[->] (b0) -- (c1); \draw[->] (b0) -- (c0);
\node[font=\scriptsize] at (0,-3.0) {année 1}; \node[font=\scriptsize] at (3,-3.0) {année 2}; \node[font=\scriptsize] at (6,-3.0) {année 3};
\node[font=\scriptsize,align=left,anchor=west] at (7.0,2.2) {payer 2.0~m\$, le prix\\revient à 1.00};
\end{tikzpicture}

\omcaption{Les trois premières années de l'arbre des prix de renouvellement~: le multiplicateur de prix du fournisseur augmente de 30~\% avec une probabilité de 0.3 à chaque renouvellement. La meilleure
politique consiste à changer de fournisseur dès que le prix a augmenté deux fois, entre l'année 3 et l'année 6~; à l'année 7, seulement après quatre hausses. Données~:
\texttt{firm.buildbuy.switch\_policy}.}\label{fig:fm:build-against-buy:tree}
\end{omfigure}

\begin{definition}[Séquestre du code source]\label{def:fm:build-against-buy:escrow}
Le \emph{séquestre du code source}\index{séquestre du code source} est un dispositif par lequel un fournisseur dépose le code source de son produit auprès d'un
agent indépendant, dont la remise au client intervient sur des événements définis, tels que l'insolvabilité du fournisseur ou l'arrêt du support du produit.
\end{definition}

Le séquestre ne supprime pas la \omterm{def:fm:build-against-buy:lockin}{dépendance au fournisseur}~; il abaisse le coût du pire cas, celui où le fournisseur disparaît, qui passe d'une reconstruction sous pression à la reprise
d'un code que l'entreprise n'a jamais exécuté. Sa valeur dépend de ce que le dépôt soit tenu à jour et qu'il puisse être compilé. Les options réelles en général (Dixit et
Pindyck, \emph{Investment under Uncertainty}, 1994) formalisent la même idée~: la flexibilité a une valeur qu'une comparaison statique ignore, et payer
pour elle (contrats courts, formats de données standard, séquestre) coûte parfois moins cher que la dépendance qu'elle prévient.

\section{Tutoriel~: sept ans d'un système de gestion des ordres}\label{tut:fm:build-against-buy:tut}

\begin{tutorial}
\textbf{Objectif.} Comparer la construction et l'achat d'un système de gestion des ordres sur sept ans, évaluer l'option de changer de fournisseur, et trouver quels paramètres
déterminent la réponse. \textbf{État final~:} les courbes de coût par horizon et le diagramme en tornade (la \cref{fig:fm:build-against-buy:tornado}).
\begin{tutsteps}
\item \textbf{Les trajectoires.} \texttt{fm\_buildbuy.paths()} énumère les coûts annuels~; \texttt{firm.buildbuy.npc} les actualise.
\item \textbf{Les seuils d'équilibre.} \texttt{breakeven\_wage} résout la \cref{prop:fm:build-against-buy:wage}~; \texttt{breakeven\_horizon} trouve l'horizon à partir duquel
construire reste moins coûteux.
\item \textbf{L'option.} \texttt{switch\_value} évalue le droit de changer de fournisseur sur l'arbre des prix de renouvellement.
\item \textbf{Le diagramme en tornade.} \texttt{fm\_buildbuy.tornado()} fait varier chaque paramètre de 20~\% à la hausse comme à la baisse et enregistre l'écart entre la construction et l'achat.
\end{tutsteps}
\end{tutorial}

\begin{omfigure}
\begin{tikzpicture}
\begin{axis}[width=10.5cm,height=5.6cm,xmin=-2,xmax=1.5,xlabel={\small construction moins achat, valeur actuelle nette des coûts sur 7 ans (millions de \$)},
  ytick={0,1,2,3,4,5,6,7},yticklabels={licence,salaire,ingénieurs de développement,ingénieurs de maintenance,ingénieurs de support,kappa,indexation,frais d'intégration},
  y dir=reverse,ymin=-0.6,ymax=7.6,tick label style={font=\footnotesize},grid=major,grid style={gray!20},table/col sep=comma,
  legend cell align=left,legend style={font=\scriptsize,at={(0.5,-0.25)},anchor=north,/tikz/every even column/.append style={column sep=8pt}},legend columns=2]
\addplot[draw=none,mark=none,error bars/.cd,x dir=both,x explicit,error bar style={line width=5pt,omDef!70},error mark=none]
  table[x=lmid,y expr=\thisrow{pos}-0.13,x error=lhalf]{figdata/desk/20-build-against-buy/tornado.csv};
\addplot[draw=none,mark=none,error bars/.cd,x dir=both,x explicit,error bar style={line width=5pt,omProp!70},error mark=none]
  table[x=hmid,y expr=\thisrow{pos}+0.13,x error=hhalf]{figdata/desk/20-build-against-buy/tornado.csv};
\draw[black,thick] (axis cs:-0.2594,-0.6) -- (axis cs:-0.2594,7.6);
\addlegendimage{omDef!70,line width=5pt}\addlegendimage{omProp!70,line width=5pt}
\legend{,,paramètre à 20~\% de moins,paramètre à 20~\% de plus}
\end{axis}
\end{tikzpicture}

\omcaption{Diagramme en tornade de la comparaison sur sept ans~: la valeur actuelle nette des coûts construction moins achat lorsque chaque paramètre varie de 20~\% à la baisse ou à la hausse~; la ligne verticale représente
le cas de base, $-0.26$ million de dollars (construction moins chère). Les valeurs négatives favorisent la construction. La licence et le salaire des ingénieurs déterminent la réponse. Données~:
\texttt{fm\_buildbuy.tornado}.}\label{fig:fm:build-against-buy:tornado}
\end{omfigure}

La licence et le salaire dominent~: une licence 20~\% moins chère rend l'achat meilleur marché de 1.17~millions de dollars que la construction, une licence 20~\% plus chère rend la construction moins coûteuse de 1.69
millions de dollars~; le salaire déplace la réponse des mêmes montants en sens inverse. Vient ensuite le nombre d'ingénieurs de développement. L'indexation
et les frais d'intégration ne pèsent guère sur sept ans. La décision repose donc sur deux chiffres qui peuvent être négociés ou mesurés~:
le montant de la licence que le fournisseur acceptera, et ce que coûtent réellement les ingénieurs de l'entreprise et ce qu'ils pourraient construire à la place.

\section{Choisir un fournisseur}

\begin{definition}[Appel d'offres (RFP)]\label{def:fm:build-against-buy:rfp}
Un \emph{appel d'offres}\index{appel d'offres} (request for proposal) est une invitation structurée adressée aux fournisseurs à proposer un produit ou un service au regard des
exigences écrites de l'entreprise, dans un format commun, afin que les propositions puissent être notées et comparées sur les fonctions, les performances, le coût, les conditions et le risque.
\end{definition}

\begin{method}[Conduire une décision de faire ou acheter]\label{met:fm:build-against-buy:run}
\begin{enumerate}
\item Rédiger les exigences, en séparant ce qui distingue l'entreprise (construire ou étendre) de ce qu'elle partage avec tous (acheter).
\item Lancer un \omterm{def:fm:build-against-buy:rfp}{appel d'offres} auprès d'au moins trois fournisseurs~; noter les fonctions, la latence et la capacité au regard des accords de niveau de service
(SLA) de l'entreprise (livre~14), et chiffrer le \omterm{def:fm:build-against-buy:tco}{coût total de possession}, et non la seule licence.
\item Estimer honnêtement le coût de construction, avec la queue de maintenance et le coût d'opportunité des ingénieurs.
\item Comparer sur la durée de vie attendue du système, en tenant compte du risque de prix et du \omterm{def:fm:build-against-buy:lockin}{coût de changement}~; évaluer l'option de changer de fournisseur et ce qu'il en coûte de la conserver.
\item Négocier des conditions qui préservent l'option~: export des données, formats standard, séquestre, plafonds de prix au renouvellement, assistance à la résiliation.
\item Réexaminer à chaque renouvellement~: la \omterm{def:fm:build-against-buy:lockin}{dépendance au fournisseur} croît avec chaque intégration.
\end{enumerate}
\end{method}

Une spécification de transmission (handoff, livre~12) entre les systèmes de l'entreprise et ceux du fournisseur est l'assurance la moins coûteuse contre la \omterm{def:fm:build-against-buy:lockin}{dépendance au fournisseur}~: si l'interface est celle de l'entreprise, un
nouveau fournisseur ou une construction interne se branche sur la même prise. Une \omterm{def:fm:technology-strategy:teams}{équipe plateforme} (chapitre 19) en a la charge.

\section{Construction~: le modèle de faire ou acheter}\label{bld:fm:build-against-buy:buildbuy}

\begin{build}
\bfield{Rôle} Le coût de construction et d'achat sur un horizon, le coût d'équilibre des ingénieurs et l'horizon d'équilibre, l'option de changer de fournisseur, et
la sensibilité de la réponse.
\bfield{Interface} \texttt{firm.buildbuy}~: \texttt{Build}, \texttt{Buy}, \texttt{build\_costs}, \texttt{buy\_costs}, \texttt{npc},
\texttt{breakeven\_wage}, \texttt{breakeven\_horizon}, \texttt{switch\_value}, \texttt{tornado}.
\bfield{Règles} Les coûts sont payés en fin d'année~; les années-ingénieur portent le coût d'opportunité sur les deux trajectoires~; la valeur de l'option n'est jamais négative.
\bfield{Tests d'acceptation} \texttt{code/firm/buildbuy/tests/}~: trajectoires de coûts sur valeurs calculées à la main~; le salaire d'équilibre égalise les coûts~; la valeur de l'option
croît avec le risque de prix~; l'ordre du diagramme en tornade.
\bfield{Pour aller plus loin} Une trajectoire mixte (noyau fournisseur avec extensions internes)~; l'internalisation comme changement~; la probabilité que la construction dépasse son budget.
\end{build}

\begin{omsources}
\item Virtu Financial, Forms 8-K du 7 novembre 2018 et du 1\up{er} mars 2019 (communiqués de presse sur l'acquisition d'Investment Technology Group).
\item A.~K.~Dixit et R.~S.~Pindyck, \emph{Investment under Uncertainty}, Princeton University Press, 1994.
\item B.~W.~Boehm, \emph{Software Engineering Economics}, Prentice-Hall, 1981.
\end{omsources}

\section{Exercices}

\begin{exercise}[$\star$]\label{exo:fm:build-against-buy:1}
Calculer le coût pour l'entreprise d'une année-ingénieur chargée, coût d'opportunité inclus, ainsi que le coût des deux premières années de construction.
\end{exercise}

\begin{exercise}[$\star$]\label{exo:fm:build-against-buy:2}
Calculer le coût de la licence la septième année, ainsi que le coût de l'achat cette année-là.
\end{exercise}

\begin{exercise}[$\star$]\label{exo:fm:build-against-buy:3}
Déterminer pour quels horizons l'achat est moins cher, et à partir de quel horizon la construction devient moins chère.
\end{exercise}

\begin{exercise}[$\star\star$]\label{exo:fm:build-against-buy:4}
Pourquoi l'actualisation favorise-t-elle l'achat dans cet exemple~?
\end{exercise}

\begin{exercise}[$\star\star$]\label{exo:fm:build-against-buy:5}
Quelle était la justification affichée par Virtu pour acheter ITG, et combien a-t-elle payé~?
\end{exercise}

\begin{exercise}[$\star\star$]\label{exo:fm:build-against-buy:6}
Pourquoi l'option de changer de fournisseur vaut-elle davantage lorsque le prix du fournisseur est plus risqué~? Qu'est-ce qui la rendrait sans valeur~?
\end{exercise}

\begin{exercise}[$\star\star\star$]\label{exo:fm:build-against-buy:7}
\emph{Programmation.} Évaluer l'option de changement avec une probabilité de hausse du prix de 0.5 au lieu de 0.3. Combien coûte l'achat en espérance, avec
et sans l'option~?
\end{exercise}

\begin{exercise}[$\star\star\star$]\label{exo:fm:build-against-buy:8}
\emph{Trouver la faille.} «~La construction ne coûte rien~: nos ingénieurs sont déjà sur la paie.~»
\end{exercise}

\section{Problème~: sept années d'un système de gestion des ordres}

\begin{problem}[{Problème du week-end --- sept années d'un système de gestion des ordres}]\label{pb:fm:build-against-buy:1}
Un directeur technique doit recommander s'il faut construire ou acheter le prochain système de gestion des ordres de la société.

\medskip\noindent\textbf{Partie I --- La décision.}
\begin{enumerate}
\item Définir \omterm{def:fm:build-against-buy:bab}{faire ou acheter} (build against buy) et le \omterm{def:fm:build-against-buy:tco}{coût total de possession}.
\item Quelles capacités une société de trading doit-elle construire, et lesquelles doit-elle acheter~?
\item Qu'a acheté Virtu en 2019, pour combien, et pourquoi~?
\item Pourquoi l'acquisition de la société est-elle le moyen le plus coûteux d'obtenir un logiciel~?
\end{enumerate}

\medskip\noindent\textbf{Partie II --- Les coûts.}
\begin{enumerate}[resume]
\item Lister les coûts annuels de chaque voie.
\item Donner les coûts actuels nets sur sept ans ainsi que les totaux non actualisés.
\item Énoncer et démontrer la \cref{prop:fm:build-against-buy:wage}.
\item Donner le salaire d'équilibre et l'horizon d'équilibre.
\end{enumerate}

\medskip\noindent\textbf{Partie III --- Verrouillage.}
\begin{enumerate}[resume]
\item Définir la \omterm{def:fm:build-against-buy:lockin}{dépendance au fournisseur} (vendor lock-in) et le \omterm{def:fm:build-against-buy:lockin}{coût de changement}.
\item Décrire l'arbre des prix de renouvellement et donner le coût espéré de l'achat avec et sans l'option de changement.
\item Définir le \omterm{def:fm:build-against-buy:escrow}{séquestre du code source} et dire ce qu'il fait et ce qu'il ne fait pas.
\item Quelles clauses contractuelles préservent l'option de changement~?
\end{enumerate}

\medskip\noindent\textbf{Partie IV --- La recommandation.}
\begin{enumerate}[resume]
\item Lire le diagramme en tornade~: quelles entrées décident~?
\item Que faudrait-il négocier avec le fournisseur~?
\item Que faudrait-il mesurer sur ses propres ingénieurs~?
\item Définir un \omterm{def:fm:build-against-buy:rfp}{appel d'offres} (RFP) et en esquisser un pour le système.
\item Comment une spécification de transmission (handoff) réduit-elle le verrouillage~?
\item Quand faut-il réexaminer la décision~?
\item Énoncer le \emph{résultat nommé}~: le coût d'une année-ingénieur chargée à partir duquel la construction
l'emporte sur l'achat sur sept ans, et la valeur de l'option de changement.
\item En deux phrases, rédiger la recommandation.
\end{enumerate}
\end{problem}

\section{Questions d'entretien}

\begin{interviewq}[$\star$ \iqroles{développeur}]\label{iq:fm:build-against-buy:1}
Votre équipe veut construire son propre gestionnaire de données de marché plutôt que d'utiliser celui d'un fournisseur. Que demanderiez-vous~?
\end{interviewq}

\begin{interviewq}[$\star$ \iqroles{trader}]\label{iq:fm:build-against-buy:2}
Pourquoi une société de trading achèterait-elle un système de gestion des ordres tout en construisant ses algorithmes d'exécution~?
\end{interviewq}

\begin{interviewq}[$\star\star$ \iqroles{chercheur}]\label{iq:fm:build-against-buy:3}
Comment évalueriez-vous l'option de changer de fournisseur~?
\end{interviewq}

\begin{interviewq}[$\star\star$ \iqroles{développeur}]\label{iq:fm:build-against-buy:4}
Qu'est-ce qui rend le produit d'un fournisseur difficile à quitter, et comment concevriez-vous pour s'en affranchir~?
\end{interviewq}

\begin{interviewq}[$\star\star$ \iqroles{risque}]\label{iq:fm:build-against-buy:5}
Quels risques un fournisseur critique ajoute-t-il, et comment les géreriez-vous~?
\end{interviewq}

\begin{interviewq}[$\star\star\star$ \iqroles{développeur, chercheur}]\label{iq:fm:build-against-buy:6}
Estimez le \omterm{def:fm:build-against-buy:tco}{coût total de possession} sur sept ans de la construction d'un système qui demande six ingénieurs par an.
\end{interviewq}
